\section{Results II: the protocol decides the ranking}
\label{sec:protocol}

The comparison of \cref{sec:results} used a probe that sees 50,000 labels. A practitioner with a 5,000-label
budget cannot run it. \Cref{tab:protocol} evaluates the same encoders at 19,400 steps with every protocol of
\cref{sec:protocols}, and \cref{tab:main} gives \lpfive and \ftfive for the other budgets.

\begin{table}[t]
\centering
\small
\caption{One set of encoders, six protocols (CIFAR-10, 10\% labels, ResNet-18, 19,400 steps; test accuracy
in \%, mean $\pm$ standard deviation over seeds). The superscript is the rank of the encoder within the
column. \simclr{}$\to$\supcon is \simclr followed by 485 steps of \supcon on the 5,000 labels
(\cref{sec:where}).}
\label{tab:protocol}
\begin{adjustbox}{max width=\linewidth}
% generated by scripts/make_assets.py -- do not edit
\begin{tabular}{lccccccc}
\toprule
 & \multicolumn{2}{c}{oracle labels} & \multicolumn{5}{c}{label budget (5,000 labels)} \\
\cmidrule(lr){2-3}\cmidrule(lr){4-8}
 & & & & \multicolumn{3}{c}{FT-5k, by recipe} & \\
\cmidrule(lr){5-7}
encoder & LP-50k & $k$NN-50k & LP-5k & default & selected & oracle & FixMatch-init \\
\midrule
\simclr & 87.8{\scriptsize$\,\pm\,$0.2}\,\textsuperscript{(3)} & 85.7{\scriptsize$\,\pm\,$0.3}\,\textsuperscript{(3)} & 85.8{\scriptsize$\,\pm\,$0.2}\,\textsuperscript{(3)} & 88.3{\scriptsize$\,\pm\,$0.5}\,\textsuperscript{(1)} & 89.4{\scriptsize$\,\pm\,$0.1}\,\textsuperscript{(2)} & 89.7{\scriptsize$\,\pm\,$0.2}\,\textsuperscript{(3)} & 94.3{\scriptsize$\,\pm\,$0.1}\,\textsuperscript{(1)} \\
BYOL & 90.7{\scriptsize$\,\pm\,$0.3}\,\textsuperscript{(1)} & 88.9{\scriptsize$\,\pm\,$0.3}\,\textsuperscript{(1)} & 88.8{\scriptsize$\,\pm\,$0.4}\,\textsuperscript{(1)} & 86.7{\scriptsize$\,\pm\,$0.3}\,\textsuperscript{(3)} & 91.3{\scriptsize$\,\pm\,$0.2}\,\textsuperscript{(1)} & 91.1{\scriptsize$\,\pm\,$0.0}\,\textsuperscript{(1)} & 92.3{\scriptsize$\,\pm\,$0.2}\,\textsuperscript{(4)} \\
\simclrsub & 73.7{\scriptsize$\,\pm\,$0.2}\,\textsuperscript{(6)} & 71.2{\scriptsize$\,\pm\,$0.3}\,\textsuperscript{(6)} & 69.2{\scriptsize$\,\pm\,$0.9}\,\textsuperscript{(6)} & 79.3{\scriptsize$\,\pm\,$0.3}\,\textsuperscript{(6)} & -- & -- & -- \\
\supcon & 83.9{\scriptsize$\,\pm\,$0.3}\,\textsuperscript{(5)} & 82.4{\scriptsize$\,\pm\,$0.1}\,\textsuperscript{(5)} & 81.5{\scriptsize$\,\pm\,$0.3}\,\textsuperscript{(4)} & 82.5{\scriptsize$\,\pm\,$0.4}\,\textsuperscript{(4)} & 82.6{\scriptsize$\,\pm\,$0.1}\,\textsuperscript{(4)} & 87.1{\scriptsize$\,\pm\,$0.4}\,\textsuperscript{(4)} & 93.9{\scriptsize$\,\pm\,$0.2}\,\textsuperscript{(2)} \\
\union & 86.2{\scriptsize$\,\pm\,$0.4}\,\textsuperscript{(4)} & 84.3{\scriptsize$\,\pm\,$0.3}\,\textsuperscript{(4)} & 80.7{\scriptsize$\,\pm\,$0.6}\,\textsuperscript{(5)} & 81.6{\scriptsize$\,\pm\,$0.2}\,\textsuperscript{(5)} & 82.1{\scriptsize$\,\pm\,$0.2}\,\textsuperscript{(5)} & 86.5{\scriptsize$\,\pm\,$0.3}\,\textsuperscript{(5)} & 93.5{\scriptsize$\,\pm\,$0.1}\,\textsuperscript{(3)} \\
\simclr{}$\to$\supcon & 89.0{\scriptsize$\,\pm\,$0.2}\,\textsuperscript{(2)} & 88.2{\scriptsize$\,\pm\,$0.0}\,\textsuperscript{(2)} & 87.7{\scriptsize$\,\pm\,$0.3}\,\textsuperscript{(2)} & 88.0{\scriptsize$\,\pm\,$0.2}\,\textsuperscript{(2)} & 89.0{\scriptsize$\,\pm\,$0.3}\,\textsuperscript{(3)} & 89.7{\scriptsize$\,\pm\,$0.1}\,\textsuperscript{(2)} & -- \\
\bottomrule
\end{tabular}

\end{adjustbox}
\end{table}

\paragraph{The ranking reverses}
At 9,700 steps the oracle-label probe orders the three main encoders as \union
(\R{union_f0.1_e100:lp50}\%) $>$ \supcon (\R{supcon_lab_f0.1_e100:lp50}\%) $\approx$ \simclr
(\R{simclr_all_e100:lp50}\%). Fine-tuning on the 5,000 labels orders them as \simclr
(\R{simclr_all_e100:ft}\%) $>$ \supcon (\R{supcon_lab_f0.1_e100:ft}\%) $>$ \union
(\R{union_f0.1_e100:ft}\%): the exact opposite. The encoder that the standard protocol puts first is last
when the labels are limited to the ones that exist, and \simclr gains
\R{cmp:simclr-union:ft:e100:diff}\pt on it (\R{cmp:simclr-union:ft:e100:ci}).

At 19,400 steps \simclr is first under every protocol, but the order of the two label-using methods still
flips: \union is \R{cmp:union-supcon:lp50:e200:diff}\pt \emph{above} \supcon with the oracle-label probe and
\R{d:union-supcon:lp5:e200}\pt \emph{below} it with \lpfive. The unlabeled images that \union adds to \supcon
improve what an oracle can read out of the features, and do not improve what the 5,000 labels can read out.
The size of the effect is also protocol-dependent: the advantage of \simclr over \supcon is
\R{cmp:simclr-supcon:lp50:e200:diff}\pt with \lpall, \R{cmp:simclr-supcon:lp5:e200:diff}\pt with \lpfive and
\R{cmp:simclr-supcon:ft:e200:diff}\pt with \ftfive.

\paragraph{Why the protocols disagree}
The label-budget protocols reuse the very labels that \supcon and \union were trained on. On their own 5,000
labeled images these encoders already separate the classes perfectly: the \lpfive probe reaches
\R{supcon_lab_f0.1_e200:lp5_train}\% training accuracy on them, against \R{simclr_all_e200:lp5_train}\% for
\simclr. A probe trained on images that the encoder has pulled into tight clusters learns a decision rule
that fits those images and transfers less well to new ones (\cref{sec:geometry}). The oracle-label probe
hides this, because 45,000 of its 50,000 training images were not labeled during pretraining.

\paragraph{Is this an artefact of the fine-tuning recipe?}
\label{sec:ftrobust}
There is a simple alternative explanation for the weak fine-tuning of the label-using encoders. Our default
recipe starts from the linear probe, which already fits the 5,000 labels perfectly; the cross-entropy
gradient is then close to zero and the encoder hardly moves. That would be a property of the recipe, not of
the encoder. We tested it with a grid of 72 recipes per encoder: head initialisation (linear probe or
random), learning rate (0.003, 0.01, 0.03, 0.1), number of passes (30, 60, 120) and a regulariser that keeps
the training loss above zero (none, label smoothing \citep{szegedy2016rethinking}, mixup
\citep{zhang2018mixup}).

\emph{Selecting the recipe with the label budget.} We first selected the recipe with the 5,000 labels only:
4,000 for training and 1,000 for validation. For \simclr this works as intended. The validation accuracy
spreads over \R{ftgrid:simclr_all_e200:valrange}\pt across the grid, the selected recipe (linear-probe head,
label smoothing) raises \ftfive from \R{simclr_all_e200:ftd}\% to \R{simclr_all_e200:ftt:pm}\%, and the
default recipe was close to the best. For \supcon and \union the selection fails for a reason that is itself
part of our finding: the 1,000 validation images were labeled images during pretraining, the encoder has
already separated them, and most recipes reach 100\% on them (the best validation score is
\R{ftgrid:supcon_lab_f0.1_e200:valmax}\% for \supcon and \R{ftgrid:union_f0.1_e200:valmax}\% for \union).
The validation score cannot rank the recipes, and the recipe it returns is no better than the default
(\R{supcon_lab_f0.1_e200:ftt}\% against \R{supcon_lab_f0.1_e200:ftd}\% for \supcon). A method that uses the
labels in pretraining therefore also uses up the only data on which a practitioner could tune what comes
next. In \cref{tab:protocol} the ``selected'' column gives, for the label-using encoders, the recipe selected
for \simclr.

\emph{An upper bound.} To see what any selection rule could achieve, we scored all 72 recipes on the test
set of seed 0, took the best one for each encoder, and applied it to the other seeds (``oracle'' column;
\cref{tab:ftsearch}). \pending{The best recipe for \supcon uses \R{ftgrid:supcon_lab_f0.1_e200:recipe} and
reaches \R{supcon_lab_f0.1_e200:fto:pm}\%, \R{d:fto-ft:supcon_lab_f0.1_e200}\pt above the default; for \union
it reaches \R{union_f0.1_e200:fto:pm}\% and for \simclr \R{simclr_all_e200:fto:pm}\%. The concern is
therefore partly justified: starting from the linear probe without regularisation under-uses the
label-using encoders. With the oracle recipe the advantage of \simclr over \supcon shrinks from
\R{cmp:simclr-supcon:ft:e200:diff}\pt to \R{cmp:simclr-supcon:fto:e200:diff}\pt
(\R{cmp:simclr-supcon:fto:e200:ci}); it does not close, and a practitioner has no label-budget way of
finding that recipe.}

\paragraph{Semi-supervised fine-tuning}
The pipeline a practitioner would use when unlabeled images are available is semi-supervised
\citep{chen2020simclrv2}. We therefore started FixMatch from each encoder. \pending{FixMatch-init reaches
\R{simclr_all_e200:fm:pm}\% from \simclr, \R{supcon_lab_f0.1_e200:fm:pm}\% from \supcon and
\R{union_f0.1_e200:fm:pm}\% from \union. The difference between \simclr and \supcon is
\R{cmp:simclr-supcon:fm:e200:sdiff}\pt (\R{cmp:simclr-supcon:fm:e200:ci}).}

\paragraph{What we take from this}
A statement such as ``method A learns better representations than method B with 10\% labels'' needs the
protocol attached. In our experiments the all-label probe is the protocol most favourable to the
label-using methods, and each protocol that respects the label budget moves the comparison towards
label-free pretraining. We recommend reporting at least one label-budget protocol next to the oracle-label
probe whenever the pretraining method itself uses labels.
